\section{Описание программной реализации}

В качестве основного языка программирования был выбран Python. При разработке вычислительных модулей использовались следующие специализированные библиотеки:
\begin{itemize}
    \item \texttt{fastecdsa} — написанная на C библиотека для работы с арифметикой эллиптических кривых над простыми полями $\mathbb{F}_p$.
    \item \texttt{SymPy} — реализация вспомогательных теоретико-числовых алгоритмов. В частности, использовались функции для факторизации целых чисел (при реализации алгоритма Похлига--Хеллмана) и вычисления систем сравнений с помощью Китайской теоремы об остатках (КТО).
\end{itemize}

Архитектура системы построена по принципу клиент-серверного взаимодействия. Вычислительное ядро отделено от интерфейса пользователя, а взаимодействие между ними организовано посредством REST API на базе фреймворка \texttt{FastAPI}.

Для удобства эксплуатации был спроектирован веб-интерфейс, позволяющий запускать решение ECDLP на заранее заготовленных кривых.

Полный исходный код проекта, включая инструкции по развертыванию, доступен в публичном репозитории: \url{https://github.com/guyder0/ecdlp}.